% TOP_PROBLEM: {"id": 4800003, "title": "Multiple ergodic averages — Problem 3", "category_id": 11, "category": {"id": 11, "name": "dynamical_systems", "display_name": "Dynamical Systems", "description": "Problems about long-term behavior of deterministic systems, Hamiltonian dynamics, and periodic orbits.", "slug": "dynamical-systems", "order_index": 11, "created_at": "2026-07-31T15:26:25.671Z"}, "status": "partially_solved"}
% FIRST_POSTED: null
% ENTRY_KIND: statement_only
% Imported from: batch15.tex; UnsolvedMath ID 4800003
\documentclass[11pt]{article}
\usepackage[margin=25mm]{geometry}
\usepackage{fontspec}
\setmainfont{Latin Modern Roman}
\usepackage{amsmath,amssymb,mathtools}
\usepackage{unicode-math}
\setmathfont{Latin Modern Math}
\usepackage{xurl,hyperref}
\hypersetup{colorlinks=true,urlcolor=blue}
\setcounter{secnumdepth}{0}
\setlength{\parindent}{0pt}
\setlength{\parskip}{5pt}
\setlength{\emergencystretch}{3em}
\newcommand{\sref}[2]{\hyperlink{source:#1:#2}{[S#2]}}
\newcommand{\cataloguescope}[1]{\paragraph{Recorded target.}#1}
\begin{document}
% UnsolvedMath ID 4800003; source catalogue code AMR-047-0003
\setcounter{section}{4800002}
\section{Multiple ergodic averages — Problem 3}\label{4800003}
{\small\sffamily UnsolvedMath ID 4800003 · Dynamical Systems}
\subsection{Definitions and mathematical statement}

\paragraph{Shared notation.}$\mathbb N=\{1,2,\ldots\}$, and $\mathbb Z,\mathbb Q,\mathbb R,\mathbb C$ denote the integers, rationals, reals and complex numbers. Any other conventions are specified locally.


A system is a probability space $(X,\mathcal X,\mu)$ with invertible measure-preserving transformations; multiple transformations are assumed to commute. Write $T^j f=f\circ T^j$, and $T^j A$ for the image of a measurable set under $T^j$. A system with one transformation is ergodic if every invariant measurable set has measure zero or one. The space $L^\infty(\mu)$ consists of essentially bounded complex measurable functions. A nilmanifold of step at most $s$ is $Y=G/\Gamma$, where $G$ is a nilpotent Lie group of step at most $s$ and $\Gamma$ is a discrete cocompact subgroup. A nilsystem is a translation $y\mapsto by$ on $Y$ with its normalized invariant Haar measure. A basic $s$-step nilsequence in $q$ variables is $\psi(n_1,\ldots,n_q)=F(b_1^{n_1}\cdots b_q^{n_q}y)$, with commuting $b_i\in G$, $y\in Y$, and continuous $F:Y\to\mathbb C$. A basic generalized nilsequence permits a bounded Riemann-integrable $F$ instead. A (generalized) nilsequence is a uniform limit of basic (generalized) nilsequences of the same step and number of variables. Fix $\ell\ge1$ and integer sequences $a_j:\mathbb N\to\mathbb Z$. They are good for $\ell$-convergence of commuting transformations when, for every system and all bounded $f_1,\ldots,f_\ell$, the averages $A_N=N^{-1}\sum_{n=1}^N\prod_jT_j^{a_j(n)}f_j$ converge in $L^2(\mu)$. They are good for $\ell$-convergence of $\ell$-step nilsystems when the same holds for every nilmanifold of step at most $\ell$ with commuting translations $T_j$ and its Haar probability measure, still for all bounded measurable functions.
\paragraph{Statement recorded in the supplied source.}
Are the following three statements equivalent for every family $(a_1,\ldots,a_\ell)$? (i) The family is good for $\ell$-convergence of arbitrary commuting transformations. (ii) It is good for $\ell$-convergence of $\ell$-step nilsystems. (iii) For every basic generalized $\ell$-step nilsequence $\psi$ in $\ell$ variables, the scalar averages $N^{-1}\sum_{n=1}^N\psi(a_1(n),\ldots,a_\ell(n))$ converge. \sref{4800003}{2}
\subsection{Short English statement}
Do convergence tests on nilsystems or on all basic generalized nilsequences characterize convergence of the corresponding averages in every commuting system?
\subsection{Sources}
\begin{description}
\item[\hypertarget{source:4800003:1}{S1}] ULAM.AI and the UnsolvedMath Contributors, \emph{UnsolvedMath: A Curated Collection of Open Mathematics Problems}, record 4800003 (AMR-047-0003). CC BY 4.0. \url{https://huggingface.co/datasets/ulamai/UnsolvedMath}.
\item[\hypertarget{source:4800003:2}{S2}] N. Frantzikinakis, Some open problems on multiple ergodic averages (2016), Sections 1.2,2.4 and3.2, Definition 3.3; Section 4.2, Problem 3, pp.2,5--6,16--17,20. \url{https://arxiv.org/pdf/1103.3808v3}.
\end{description}
\label{end:4800003}
\end{document}
